% Options for packages loaded elsewhere
\PassOptionsToPackage{unicode}{hyperref}
\PassOptionsToPackage{hyphens}{url}
\PassOptionsToPackage{dvipsnames,svgnames,x11names}{xcolor}
%
\documentclass[
  10pt,
]{article}
\usepackage{amsmath,amssymb}
\usepackage{iftex}
\ifPDFTeX
  \usepackage[T1]{fontenc}
  \usepackage[utf8]{inputenc}
  \usepackage{textcomp} % provide euro and other symbols
\else % if luatex or xetex
  \usepackage{unicode-math} % this also loads fontspec
  \defaultfontfeatures{Scale=MatchLowercase}
  \defaultfontfeatures[\rmfamily]{Ligatures=TeX,Scale=1}
\fi
\usepackage{lmodern}
\ifPDFTeX\else
  % xetex/luatex font selection
\fi
% Use upquote if available, for straight quotes in verbatim environments
\IfFileExists{upquote.sty}{\usepackage{upquote}}{}
\IfFileExists{microtype.sty}{% use microtype if available
  \usepackage[]{microtype}
  \UseMicrotypeSet[protrusion]{basicmath} % disable protrusion for tt fonts
}{}
\makeatletter
\@ifundefined{KOMAClassName}{% if non-KOMA class
  \IfFileExists{parskip.sty}{%
    \usepackage{parskip}
  }{% else
    \setlength{\parindent}{0pt}
    \setlength{\parskip}{6pt plus 2pt minus 1pt}}
}{% if KOMA class
  \KOMAoptions{parskip=half}}
\makeatother
\usepackage{xcolor}
\usepackage[margin=2.3cm]{geometry}
\usepackage{longtable,booktabs,array}
\usepackage{calc} % for calculating minipage widths
% Correct order of tables after \paragraph or \subparagraph
\usepackage{etoolbox}
\makeatletter
\patchcmd\longtable{\par}{\if@noskipsec\mbox{}\fi\par}{}{}
\makeatother
% Allow footnotes in longtable head/foot
\IfFileExists{footnotehyper.sty}{\usepackage{footnotehyper}}{\usepackage{footnote}}
\makesavenoteenv{longtable}
\setlength{\emergencystretch}{3em} % prevent overfull lines
\providecommand{\tightlist}{%
  \setlength{\itemsep}{0pt}\setlength{\parskip}{0pt}}
\setcounter{secnumdepth}{-\maxdimen} % remove section numbering
\newlength{\cslhangindent}
\setlength{\cslhangindent}{1.5em}
\newlength{\csllabelwidth}
\setlength{\csllabelwidth}{3em}
\newlength{\cslentryspacingunit} % times entry-spacing
\setlength{\cslentryspacingunit}{\parskip}
\newenvironment{CSLReferences}[2] % #1 hanging-ident, #2 entry spacing
 {% don't indent paragraphs
  \setlength{\parindent}{0pt}
  % turn on hanging indent if param 1 is 1
  \ifodd #1
  \let\oldpar\par
  \def\par{\hangindent=\cslhangindent\oldpar}
  \fi
  % set entry spacing
  \setlength{\parskip}{#2\cslentryspacingunit}
 }%
 {}
\usepackage{calc}
\newcommand{\CSLBlock}[1]{#1\hfill\break}
\newcommand{\CSLLeftMargin}[1]{\parbox[t]{\csllabelwidth}{#1}}
\newcommand{\CSLRightInline}[1]{\parbox[t]{\linewidth - \csllabelwidth}{#1}\break}
\newcommand{\CSLIndent}[1]{\hspace{\cslhangindent}#1}
\ifLuaTeX
\usepackage[bidi=basic]{babel}
\else
\usepackage[bidi=default]{babel}
\fi
\babelprovide[main,import]{spanish}
% get rid of language-specific shorthands (see #6817):
\let\LanguageShortHands\languageshorthands
\def\languageshorthands#1{}
\ifLuaTeX
  \usepackage{selnolig}  % disable illegal ligatures
\fi
\IfFileExists{bookmark.sty}{\usepackage{bookmark}}{\usepackage{hyperref}}
\IfFileExists{xurl.sty}{\usepackage{xurl}}{} % add URL line breaks if available
\urlstyle{same}
\hypersetup{
  pdftitle={Publicaciones previstas --- SedCoreTool},
  pdfauthor={Santiago Casanova-Arenillas},
  pdflang={es},
  colorlinks=true,
  linkcolor={Maroon},
  filecolor={Maroon},
  citecolor={Blue},
  urlcolor={Blue},
  pdfcreator={LaTeX via pandoc}}

\title{Publicaciones previstas --- SedCoreTool}
\usepackage{etoolbox}
\makeatletter
\providecommand{\subtitle}[1]{% add subtitle to \maketitle
  \apptocmd{\@title}{\par {\large #1 \par}}{}{}
}
\makeatother
\subtitle{Cartera, métodos, defensa ante revisores y borradores}
\author{Santiago Casanova-Arenillas}
\date{07-11-2026 · software}

\begin{document}
\maketitle

{
\hypersetup{linkcolor=}
\setcounter{tocdepth}{2}
\tableofcontents
}
\hypertarget{publicaciones-previstas-sedcoretool-antes-sedcoretool-e-ichantool}{%
\section{Publicaciones previstas --- SedCoreTool (antes SedCoreTool e
IchAnTool)}\label{publicaciones-previstas-sedcoretool-antes-sedcoretool-e-ichantool}}

Borradores, revistas candidatas y estado de cada trabajo. Estado al
\textbf{01-11-2026} (software en la entrega V28). Todo lo marcado
\texttt{{[}TODO{]}} está pendiente; ningún resultado de estos borradores
es definitivo.

\begin{quote}
Aviso: si el repositorio es \textbf{público}, no subir aquí datos ni
resultados inéditos privados hasta que el manuscrito esté enviado o
publicado como preprint (o pasar los borradores a un repositorio
privado).
\end{quote}

\hypertarget{cartera-de-publicaciones}{%
\subsection{Cartera de publicaciones}\label{cartera-de-publicaciones}}

\begin{longtable}[]{@{}
  >{\raggedright\arraybackslash}p{(\columnwidth - 8\tabcolsep) * \real{0.0222}}
  >{\raggedright\arraybackslash}p{(\columnwidth - 8\tabcolsep) * \real{0.3333}}
  >{\raggedright\arraybackslash}p{(\columnwidth - 8\tabcolsep) * \real{0.3333}}
  >{\raggedright\arraybackslash}p{(\columnwidth - 8\tabcolsep) * \real{0.1778}}
  >{\raggedright\arraybackslash}p{(\columnwidth - 8\tabcolsep) * \real{0.1333}}@{}}
\toprule\noalign{}
\begin{minipage}[b]{\linewidth}\raggedright
\#
\end{minipage} & \begin{minipage}[b]{\linewidth}\raggedright
Trabajo y aporte
\end{minipage} & \begin{minipage}[b]{\linewidth}\raggedright
Revista principal (alternativas)
\end{minipage} & \begin{minipage}[b]{\linewidth}\raggedright
Depende de
\end{minipage} & \begin{minipage}[b]{\linewidth}\raggedright
Estado
\end{minipage} \\
\midrule\noalign{}
\endhead
\bottomrule\noalign{}
\endlastfoot
1 & \textbf{Software: SedCoreTool.} Herramienta abierta y reproducible &
JOSS (Applied Computing and Geosciences; Computers \& Geosciences) &
Release con DOI & Borrador en inglés \\
2 & \textbf{Dataset anotado y segmentación automática.} Datos abiertos y
método & Earth System Science Data (Scientific Data; AI in Geosciences;
Computers \& Geosciences) & 1 e imágenes con permiso & Esquema \\
3 & \textbf{Validación frente a índices visuales y aplicación
paleoambiental.} Interpretación geológica & Palaeogeography,
Palaeoclimatology, Palaeoecology (Ichnos; Marine Geology; Sedimentary
Geology) & 1 (y 2) & Esquema \\
4 & \textbf{Clasificación automática de intervalos con datos
multiproxy.} Método de aprendizaje automático & Computers \& Geosciences
(AI in Geosciences; Applied Computing and Geosciences; Marine Geology) &
1 y sondeos con eventos descritos & Esquema y primeras pruebas \\
5 & \textbf{Nota breve en español.} Difusión en España & Geogaceta
(Spanish Journal of Palaeontology) & 1 & Idea \\
\end{longtable}

\textbf{Orden recomendado:} 1 $\rightarrow$ 5 $\rightarrow$ 4 $\rightarrow$ 2 $\rightarrow$ 3.

\begin{itemize}
\tightlist
\item
  El 1 da el DOI que citan los demás; el 5 se puede escribir a la vez
  con el mismo material.
\item
  El 4 es el más avanzado en software (V26--V27: familias,
  clasificadores, red neuronal, validación) y solo necesita datos con
  eventos descritos de forma independiente.
\item
  El 2 y el 3 dependen de trabajo con otras personas (anotadores,
  observadores a ciegas) y llevan más tiempo.
\end{itemize}

\hypertarget{quuxe9-puede-afirmar-hoy-cada-uno-resumen}{%
\subsection{Qué puede afirmar hoy cada uno
(resumen)}\label{quuxe9-puede-afirmar-hoy-cada-uno-resumen}}

\begin{longtable}[]{@{}
  >{\raggedright\arraybackslash}p{(\columnwidth - 4\tabcolsep) * \real{0.0270}}
  >{\raggedright\arraybackslash}p{(\columnwidth - 4\tabcolsep) * \real{0.4865}}
  >{\raggedright\arraybackslash}p{(\columnwidth - 4\tabcolsep) * \real{0.4865}}@{}}
\toprule\noalign{}
\begin{minipage}[b]{\linewidth}\raggedright
\#
\end{minipage} & \begin{minipage}[b]{\linewidth}\raggedright
Se puede afirmar
\end{minipage} & \begin{minipage}[b]{\linewidth}\raggedright
No se puede afirmar todavía
\end{minipage} \\
\midrule\noalign{}
\endhead
\bottomrule\noalign{}
\endlastfoot
1 & Reproduce exactamente los resultados de 2020; es determinista;
integra imágenes, registros, intervalos, clasificación, mapas y figuras
& Rendimiento de la detección automática en datos reales \\
2 & Las herramientas de anotación existen y producen máscaras abiertas;
flujo U-Net completo (exportar, entrenar en el programa o en Python,
evaluar, usar); rasgos DINOv2 para icnotaxones (V0.15--V0.17) &
Precisión de la segmentación y de la identificación (faltan datos
anotados) \\
3 & La app registra BI/ii manuales junto a los automáticos &
Concordancia con observadores (falta el experimento) \\
4 & La cadena funciona con verdad conocida (sintético: RF $\approx$ 91 \%, red $\approx$
80 \%) & Capacidad predictiva real: con los datos privados del autor, 18
\% de acierto equilibrado dejando un sondeo fuera (7 clases) \\
5 & Funcionalidad y ejemplo con datos públicos & --- \\
\end{longtable}

\hypertarget{documentaciuxf3n-de-muxe9todos}{%
\subsection{Documentación de
métodos}\label{documentaciuxf3n-de-muxe9todos}}

Explicación sencilla del aprendizaje automático (segmentación U-Net,
icnotaxones, validación): GUIA\_APRENDIZAJE\_AUTOMATICO.md. Qué
algoritmos usa cada artículo, qué se puede afirmar, qué falta validar y
respuestas preparadas a revisores: METODOS\_Y\_DEFENSA.md. Detalle
técnico con fórmulas: ../docs/TECHNICAL\_DOCUMENTATION.md. Revistas:
journals.md.

\hypertarget{antes-de-enviar}{%
\subsection{Antes de enviar}\label{antes-de-enviar}}

\begin{itemize}
\tightlist
\item
  Confirmar coautorías (F.J. Rodríguez-Tovar, F. Martínez-Ruiz, J.
  Dorador u otros según aporten datos o validación).
\item
  ORCID de todos los autores.
\item
  Release con DOI de Zenodo del código citado (indicar la entrega,
  p.~ej. V28).
\item
  Comprobar en la web de cada revista: alcance, formato, APC y política
  de preprints.
\item
  Preprint opcional en EarthArXiv.
\end{itemize}

\newpage

\hypertarget{muxe9todos-y-defensa-ante-revisores}{%
\section{Métodos y defensa ante
revisores}\label{muxe9todos-y-defensa-ante-revisores}}

Estado al 01-11-2026 (software en la entrega \textbf{V28}, aplicación
web 2.0.1). Cartera completa en README.md.

\hypertarget{lo-que-hay-que-saber-del-software-antes-de-escribir-cualquier-artuxedculo}{%
\subsection{Lo que hay que saber del software antes de escribir
cualquier
artículo}\label{lo-que-hay-que-saber-del-software-antes-de-escribir-cualquier-artuxedculo}}

\begin{itemize}
\tightlist
\item
  \textbf{Nombre:} desde V0.52 el programa se llama \textbf{SedCoreTool}
  (antes SedCoreTool, V16--V0.51, e IchAnTool). En los artículos:
  «SedCoreTool (formerly IchAnTool; Casanova-Arenillas et al., 2020)».
  El código del artículo de 2020 sigue en Zenodo
  (doi:10.5281/zenodo.3518837), en \texttt{legacy/} y en la etiqueta
  \texttt{v2.1}. Citar la entrega usada (p.~ej. «V28») y, al enviar, la
  release con DOI.
\item
  \textbf{Icnotaxones (imágenes):} tres modelos (sondeo, proyecto, todos
  los proyectos; §10.6 bis) y tres métodos (k-NN, bosque aleatorio,
  Bayes ingenuo; §10.6 ter). Desde V22 la evaluación agrupa por imagen
  (sección) y cuenta como error cualquier predicción de un taxón ausente
  del conjunto de prueba. Hay que decir qué modelo y método se usó y su
  acierto equilibrado. Las «probabilidades» de RF, NB y k-NN son
  proporciones de voto o verosimilitudes normalizadas, no probabilidades
  calibradas.
\item
  \textbf{Multiproxy (V26--V27):} familias de intervalos (p.~ej. periodo
  frío/cálido; baja oxigenación con sapropel/ORL) que se clasifican a
  mano en cualquier sondeo, y \textbf{clasificadores definidos por el
  usuario} (qué familias, con qué variables y familias de entrada, qué
  modelo: bosque aleatorio, red neuronal MLP, k-NN, Bayes, k-medias),
  validados dejando un sondeo fuera (§18.4--§18.8). Solo se puntúan
  intervalos descritos, no los etiquetados por reglas. Desde V0.37 las
  familias madre e hijas de la familia objetivo se excluyen de las
  entradas y el escalado global se ajusta en cada pliegue con el
  entrenamiento.
\item
  \textbf{Icnotaxones (V0.37):} la evaluación deja fuera secciones (o
  sondeos, a elegir) enteras: pliegues agrupados con $\geq$ 5 grupos, un
  grupo cada vez con 2--4 y, con uno solo, trazas sueltas con aviso de
  acierto muy optimista. El acierto mostrado es el de la evaluación que
  elige el método, así que es optimista hasta medirlo en un sondeo
  reservado.
\item
  \textbf{Prueba final (V0.49):} los sondeos marcados como reservados no
  entrenan ni intervienen al elegir el modelo (en icnotaxones ni en
  multiproxy) y se usan una sola vez para el acierto final; cada
  repetición queda registrada con fecha y modelo, de modo que el
  artículo puede declarar cuántas veces se miró el conjunto reservado.
\item
  \textbf{Selección de modelo:} comparar varios clasificadores con la
  misma validación y quedarse con el mejor \textbf{sesga el acierto al
  alza}. Un artículo debe informar de todos los modelos probados o
  reservar sondeos para una prueba final independiente (pendiente de
  implementar en la interfaz).
\item
  \textbf{Sondeos compartidos entre proyectos (V23):} cuentan una sola
  vez en mapas y análisis entre proyectos, y los ejemplos de icnotaxones
  repetidos se descartan. La coincidencia es por nombre: revisar los
  marcados «en N proyectos».
\item
  \textbf{Mapas:} los valores usan solo datos reales; la interpolación
  (IDW, kriging, TIN, RBF\ldots; §19.5) mide distancias en km y es
  exploratoria. Los mapas de z por sondeo eliminan el gradiente absoluto
  entre sondeos: úsense para patrones relativos. Los fondos (Esri,
  GEBCO) son contexto visual y se citan en el pie.
\item
  \textbf{Modo identificación (V32):} registra presencia de icnotaxones
  (marcas puntuales), no abundancia: con él no se pueden dar
  porcentajes, BI ni ii, solo distribución de presencias por intervalo.
  Indicar en los métodos qué tramos se analizaron en este modo.
\item
  \textbf{Figuras de secciones (V32):} paneles original / tratada /
  trazas resaltadas / flechas, al estilo de Dorador y Rodríguez-Tovar;
  las trazas resaltadas son las clasificadas en el programa (incluidas
  las predicciones aceptadas si así está configurado), lo que debe
  constar en el pie.
\item
  \textbf{Datos editados a mano (V0.35):} los valores cambiados en la
  hoja de Datos se guardan con el original y la fecha y se listan en el
  CSV; hay que declararlos si se usan. Las columnas calculadas
  interpolan linealmente las columnas usadas a cada profundidad cuando
  se midieron a profundidades distintas.
\item
  \textbf{Modelo de edad (V0.28):} se puede importar un modelo de Bacon
  (mediana o media e intervalo del 95 \%); el programa no ejecuta Bacon
  ni propaga las iteraciones del MCMC. Las edades por intervalo se
  interpolan linealmente sobre la tabla de Bacon; citar Bacon (Blaauw y
  Christen, 2011) y la curva de calibración usada.
\item
  \textbf{Elección del realce (V0.28):} con trazas identificadas, el
  criterio es la separación (2·AUC − 1) entre esas trazas y la matriz;
  si las mismas trazas se usan después para evaluar la detección, la
  evaluación es optimista: elegir el ajuste con unas secciones y evaluar
  con otras. Sin trazas identificadas (V0.37) la puntuación es el
  cociente de Fisher en el umbral de Otsu multiplicado por la coherencia
  espacial κ de la separación, de modo que el ruido sin estructura
  puntúa $\approx$ 0 con cualquier realce; no es un método publicado y se
  presenta como propuesta inicial.
\item
  \textbf{Trazas solapadas (V0.20):} una traza puede estar dentro de
  otra. La bioturbación total cuenta el solape una vez ($\leq$ 100 \%), pero
  cada icnotaxón cuenta su superficie completa, así que la suma de los
  porcentajes por icnotaxón puede superar al total (y a 100 \%). Hay que
  decirlo en los métodos y no presentar los porcentajes por icnotaxón
  como partes de un todo. Antes de V0.20 cada píxel tenía una sola traza
  (la última dibujada).
\item
  \textbf{Datos privados del autor:} LC06, LC14 y LC20 tienen
  profundidad supuesta (secciones apiladas, §21.8 ter): no usarlos en
  una figura publicada sin verificarla. Las cifras obtenidas con
  versiones anteriores a V22 deben recalcularse.
\end{itemize}

Para cada artículo previsto: qué partes de IchAnTool usa, qué se puede
afirmar con lo que ya está hecho, qué falta comprobar antes de enviarlo
y las preguntas que probablemente hará un revisor, con una respuesta
preparada.

Los detalles de cada algoritmo (fórmulas, parámetros, archivos de
código) están en \texttt{docs/TECHNICAL\_DOCUMENTATION.md}. Aquí se cita
la sección (§). El historial de cambios está en
\texttt{docs/HISTORIAL\_VERSIONES.md}.

Regla general: \textbf{solo se afirma lo que está medido}. Lo que no
está validado se presenta como limitación o trabajo futuro.

\begin{center}\rule{0.5\linewidth}{0.5pt}\end{center}

\hypertarget{artuxedculo-1-software-sedcoretool-joss}{%
\subsection{Artículo 1 --- Software: SedCoreTool
(JOSS)}\label{artuxedculo-1-software-sedcoretool-joss}}

\hypertarget{quuxe9-describe}{%
\subsubsection{Qué describe}\label{quuxe9-describe}}

El paquete Python y la aplicación web como herramienta abierta de
análisis de sondeos: series continuas de bioturbación a partir de
imágenes, junto con registros multiproxy, intervalos descritos,
clasificación automática de intervalos, mapas por edad y figuras,
organizado en proyectos.

\hypertarget{muxe9todos-que-usa}{%
\subsubsection{Métodos que usa}\label{muxe9todos-que-usa}}

\begin{longtable}[]{@{}
  >{\raggedright\arraybackslash}p{(\columnwidth - 2\tabcolsep) * \real{0.7273}}
  >{\raggedright\arraybackslash}p{(\columnwidth - 2\tabcolsep) * \real{0.2727}}@{}}
\toprule\noalign{}
\begin{minipage}[b]{\linewidth}\raggedright
Método
\end{minipage} & \begin{minipage}[b]{\linewidth}\raggedright
§
\end{minipage} \\
\midrule\noalign{}
\endhead
\bottomrule\noalign{}
\endlastfoot
Escala de profundidad y zona de análisis & §2 \\
Porcentaje por intervalos, multiclase y exclusiones & §3, §10 \\
BI e ii (tramos de la V2, configurables) & §4 \\
Modelo de edad (interpolación lineal) & §5 \\
Realce, detección y edición (aplicación web) & §7--§10 \\
Identificación asistida k-NN (aplicación web) & §12 \\
Formatos de exportación & §24 \\
Resolución de trabajo, proyectos & §6, §22 \\
Registro compuesto y registros externos & §15--§16 \\
Herramienta multiproxy: familias, clasificadores (RF, MLP, k-NN, NB),
validación & §18--§18.8 \\
Mapas: proyecciones, interpolación (IDW, kriging, TIN, RBF\ldots),
exportación & §19 \\
Gráficos de datos (tipo Grapher) y hoja de cálculo & §17.5 \\
Organización por proyectos y sondeos, análisis entre proyectos, figuras
& §21 \\
Proyecto de entrenamiento y recorrido guiado & §30 \\
\end{longtable}

\hypertarget{quuxe9-se-puede-afirmar-hoy}{%
\subsubsection{Qué se puede afirmar
hoy}\label{quuxe9-se-puede-afirmar-hoy}}

\begin{itemize}
\tightlist
\item
  Integra en una sola herramienta las imágenes, la bioturbación y otros
  registros del mismo sondeo, en profundidad o en edad, con exportación
  vectorial (§15). Demostrado con datos públicos del IODP (sondeo de
  prueba U1385I-1H).
\item
  Reproduce exactamente los resultados publicados en 2020 (prueba de
  regresión, §31.1).
\item
  La aplicación web da los mismos BI e ii que la versión publicada en la
  imagen de ejemplo (100 \%). Con la misma máscara la reproduce
  exactamente; con el decodificador JPEG del navegador el porcentaje
  difiere $\leq$ 0,56 puntos (media 0,13) por cómo se reconstruyen los bordes
  de las manchas, no por los pesos de gris (corregido en V0.49, §3.4,
  §13.2). Desde V0.49 las dos comparaciones son pruebas automáticas.
\item
  Es determinista y el archivo de proyecto permite reproducir cualquier
  análisis (§33).
\item
  Tiene pruebas automáticas en integración continua (§31.1).
\item
  Integra en el mismo proyecto la interpretación icnológica, registros
  numéricos de cualquier tipo, intervalos descritos jerárquicos y una
  clasificación automática de intervalos (§18). Sobre esta parte solo
  puede afirmarse que funciona como herramienta; su capacidad de
  detectar eventos reales no está validada (ver artículo 4).
\end{itemize}

\hypertarget{quuxe9-falta-antes-de-enviarlo}{%
\subsubsection{Qué falta antes de
enviarlo}\label{quuxe9-falta-antes-de-enviarlo}}

\begin{itemize}
\tightlist
\item[$\square$]
  Publicar en PyPI y crear la release con DOI en Zenodo.
\item[$\square$]
  Documentación de uso (tutorial con un testigo real) y notebook de
  ejemplo.
\item[$\square$]
  Capa de exclusión también en el paquete Python, para igualar las dos
  implementaciones.
\item[$\square$]
  Alojar la aplicación en GitHub Pages.
\item[$\square$]
  Decidir la licencia definitiva.
\item[$\square$]
  Revisar la bibliografía.
\item[$\square$]
  Dividir \texttt{app.js} en módulos (los revisores de JOSS miran la
  calidad del código).
\item[$\square$]
  Decidir el alcance: un único artículo de software con todo, o dejar la
  parte multiproxy para el artículo 4.
\end{itemize}

\hypertarget{preguntas-probables-de-revisores-joss-revisa-tambiuxe9n-el-cuxf3digo}{%
\subsubsection{Preguntas probables de revisores (JOSS revisa también el
código)}\label{preguntas-probables-de-revisores-joss-revisa-tambiuxe9n-el-cuxf3digo}}

\textbf{¿Qué aporta frente a ImageJ/Fiji, que ya cuenta píxeles?} ImageJ
da un área total. IchAnTool da una \textbf{serie continua por
profundidad o edad} con BI, ii y porcentaje por icnotaxón, lista para
correlacionar con datos geoquímicos. Gestiona varias secciones, zonas
excluidas y modelo de edad en un solo flujo. Ese flujo en ImageJ
requiere macros propias.

\textbf{¿Cómo sé que los resultados son correctos?} Pruebas con casos
calculados a mano y prueba de regresión frente a los resultados
publicados (§31.1). Ambas se ejecutan automáticamente en cada cambio.

\textbf{¿Cómo se eligen los parámetros de realce?} Con los ajustes
publicados (Dorador et al., 2014; Miguez-Salas et al., 2019) o con la
propuesta automática (§7: separación de las trazas identificadas o, sin
ellas, Fisher × coherencia espacial; antes de V0.28, separabilidad de
Otsu), que es propia y se presenta como punto de partida. No hay en la
bibliografía un método publicado de preselección automática (revisión
V12).

\textbf{¿Los ajustes «de artículo» reproducen fielmente los protocolos
publicados?} Reproducen la secuencia y los parámetros publicados donde
existen (§7). Las diferencias están documentadas: los niveles se
calculan por percentiles porque el artículo los elige en cada imagen; la
gamma es automática dentro del rango publicado; el contraste y el brillo
de Photoshop solo tienen equivalentes aproximados; y el CLAHE es por
mosaicos (Zuiderveld, 1994) frente a la ventana deslizante de Fiji. Los
parámetros efectivos de cada sección se guardan en el proyecto.

\textbf{¿Por qué dos implementaciones (Python y JavaScript)?} Python
para análisis por lotes y reproducibilidad en scripts. Web para el
trabajo interactivo sin instalar nada. Comparten las mismas definiciones
y las pruebas comprueban los mismos casos en las dos (§1, §31.1).

\textbf{¿Por qué BI e ii salen solo del porcentaje?} Es la componente
cuantificable de esas escalas. Se documenta que no sustituye la
descripción completa del icnofábrica, y la aplicación guarda el índice
manual del usuario (§4.3).

\begin{center}\rule{0.5\linewidth}{0.5pt}\end{center}

\hypertarget{artuxedculo-2-conjunto-de-datos-anotado-y-segmentaciuxf3n-automuxe1tica}{%
\subsection{Artículo 2 --- Conjunto de datos anotado y segmentación
automática}\label{artuxedculo-2-conjunto-de-datos-anotado-y-segmentaciuxf3n-automuxe1tica}}

\hypertarget{quuxe9-describiruxe1}{%
\subsubsection{Qué describirá}\label{quuxe9-describiruxe1}}

Un conjunto abierto de imágenes de testigos con máscaras de bioturbación
(y de icnotaxones), anotado con la aplicación web, y la evaluación de
métodos automáticos de detección e identificación.

\hypertarget{muxe9todos-que-usa-1}{%
\subsubsection{Métodos que usa}\label{muxe9todos-que-usa-1}}

\begin{longtable}[]{@{}
  >{\raggedright\arraybackslash}p{(\columnwidth - 2\tabcolsep) * \real{0.7273}}
  >{\raggedright\arraybackslash}p{(\columnwidth - 2\tabcolsep) * \real{0.2727}}@{}}
\toprule\noalign{}
\begin{minipage}[b]{\linewidth}\raggedright
Método
\end{minipage} & \begin{minipage}[b]{\linewidth}\raggedright
§
\end{minipage} \\
\midrule\noalign{}
\endhead
\bottomrule\noalign{}
\endlastfoot
Herramientas de anotación (varita, pincel, lazo, exclusión) & §10 \\
Detección clásica: umbral adaptativo + morfología (método base) & §8 \\
Descriptores de objeto & §11 \\
k-NN con rechazo de novedad (método base de identificación) & §12 \\
Rasgos de un modelo de visión preentrenado (DINOv2) como entrada del
clasificador (\textbf{V0.17, sin evaluar con datos reales}) & §12.6 \\
Entrenamiento de la U-Net en el programa de escritorio (mismo código que
\texttt{sedcoretool\ unet}) & §9 \\
Segmentación aprendida U-Net: exportación de datos, entrenamiento,
evaluación y uso en la aplicación (\textbf{implementada en V0.15, sin
entrenar con datos reales}) & §9 \\
Formato de máscara y paleta & §24 \\
\end{longtable}

\hypertarget{quuxe9-se-puede-afirmar-hoy-1}{%
\subsubsection{Qué se puede afirmar
hoy}\label{quuxe9-se-puede-afirmar-hoy-1}}

Nada sobre la precisión en datos reales: todavía no hay datos anotados.
Solo que las herramientas existen y producen máscaras en un formato
abierto.

\begin{itemize}
\item
  Desde V0.15 el flujo U-Net completo existe y está probado con datos
  sintéticos: exportación de imágenes y máscaras con píxeles «ignorar»
  (zonas excluidas y trazas sin revisar), teselado y división por
  sondeo, entrenamiento, evaluación sobre imágenes completas de sondeos
  independientes y ejecución en la aplicación. Las reglas de teselado,
  normalización, unión y reescalado son idénticas en Python y en la
  aplicación (prueba de paridad). \textbf{No hay ningún resultado de
  segmentación con testigos reales}: no se puede citar ningún Dice ni
  IoU hasta entrenar y evaluar.
\item
  Desde V0.17 el clasificador de icnotaxones puede usar, en vez de los
  12 descriptores, el vector de 384 valores de un modelo de visión
  preentrenado (DINOv2). Se puede afirmar que la opción existe y cómo se
  calcula; \textbf{no} que mejore la identificación hasta compararla en
  testigos independientes. Desde V0.45 la descarga predeterminada fija
  el commit de DINOv2 y el SHA-256 de sus pesos ONNX (documentación
  técnica §12.6); esto identifica los pesos para reproducibilidad, sin
  aportar una evaluación de su precisión. En V0.48 se comprobó la
  inferencia del modelo real con ONNX Runtime Web dentro de Electron y
  una entrada sintética: 384 valores finitos, norma del vector
  redondeado $\approx$ 0,999998555, SHA-256 coincidente. Es una comprobación
  técnica de ejecución, sin medida de precisión ni comparación con
  testigos reales.
\item
  Desde V0.17 la U-Net se puede entrenar desde el programa de
  escritorio. Es el mismo código que la línea de órdenes (mismas
  teselas, pérdida, aumentos y evaluación), así que los resultados son
  comparables; los modelos guardan sondeos de cada conjunto, semilla e
  historial.
\item
  Decisiones defendibles ya fijadas en el código: entrada sin realce,
  resolución común en px/cm, sin giros de 90° en los aumentos, trazas
  sin revisar ignoradas (no tratadas como matriz), métricas sobre
  imágenes completas y no sobre teselas.
\item
  La identificación asistida de icnotaxones ya tiene evaluación por
  imagen (V22); falta medirla en testigos independientes anotados por
  expertos.
\end{itemize}

\hypertarget{experimentos-necesarios}{%
\subsubsection{Experimentos necesarios}\label{experimentos-necesarios}}

\begin{enumerate}
\def\labelenumi{\arabic{enumi}.}
\tightlist
\item
  \textbf{Datos:} testigos con imágenes de licencia clara (IODP/ODP son
  públicas; las imágenes tratadas del grupo de la UGR, con permiso).
  Anotar la resolución y el equipo de cada imagen.
\item
  \textbf{Protocolo de anotación escrito}, con criterios de qué es
  traza, qué es artefacto y qué se excluye.
\item
  \textbf{Concordancia entre anotadores:} al menos dos expertos anotan
  el mismo subconjunto. Se mide IoU (intersección sobre unión) y la
  diferencia de porcentaje por intervalo.
\item
  \textbf{División por testigo}, no por recortes: los testigos de prueba
  no se usan para entrenar ni para ajustar parámetros. Evita resultados
  inflados por autocorrelación espacial.
\item
  \textbf{Métricas:}

  \begin{itemize}
  \tightlist
  \item
    IoU y F1 por píxel.
  \item
    Error absoluto del porcentaje por intervalo.
  \item
    Coincidencia de BI e ii.
  \item
    Para la identificación: matriz de confusión y F1 por clase.
  \end{itemize}
\item
  \textbf{Comparaciones:}

  \begin{itemize}
  \tightlist
  \item
    umbral adaptativo (método base);
  \item
    flujo Fiji/ICY de Miguez-Salas et al.~(2019): CLAHE + umbral o
    K-medias + rellenar huecos, ya implementado;
  \item
    tratamiento de Dorador et al.~(2014), ya implementado;
  \item
    U-Net, comparable con Kikuchi y Naruse (2024, Dice 0,892 para
    \emph{Phycosiphon});
  \item
    concordancia entre humanos, que es el techo razonable.
  \end{itemize}
\end{enumerate}

\hypertarget{preguntas-probables-de-revisores}{%
\subsubsection{Preguntas probables de
revisores}\label{preguntas-probables-de-revisores}}

\textbf{¿Por qué un umbral adaptativo como base y no directamente deep
learning?} Porque es reproducible, no necesita datos de entrenamiento y
marca el nivel mínimo que un método aprendido debe superar. Se compara
con U-Net en el mismo conjunto de prueba.

\textbf{¿Por qué se ignoran las trazas no revisadas en vez de contarlas
como matriz?} Porque una traza que el experto no ha confirmado no es un
negativo fiable. Tratarla como matriz enseñaría a la red a no detectar
trazas reales; ignorarla solo reduce los píxeles de entrenamiento. Las
clases que no cuentan (artefactos) sí son matriz: el experto las ha
revisado.

\textbf{¿Por qué no se gira la imagen 90° para aumentar datos?} La
orientación de las trazas respecto a la estratificación (vertical,
horizontal, oblicua) es un criterio de identificación. Solo se voltea
(horizontal y vertical) y se varían brillo, contraste y gamma.

\textbf{¿Por qué usar un modelo de visión entrenado con fotos generales
y no uno propio?} Porque con pocos ejemplos por icnotaxón no se puede
entrenar una red de clasificación fiable. Los modelos autosupervisados
como DINOv2 dan descriptores genéricos de forma y textura que funcionan
bien con k-NN y pocos ejemplos. La red no se ajusta a nuestros datos (no
hay riesgo de sobreajuste en ella); solo el k-NN aprende. Se compara con
los 12 descriptores en el mismo conjunto de prueba; si no mejora, se
dice.

\textbf{¿Un modelo reanudado tras una interrupción es equivalente?}
Continúa desde el último punto guardado (red, optimizador, época, mejor
resultado). No es idéntico bit a bit a uno sin interrupción (el orden
aleatorio de los aumentos no se guarda). Para el artículo se entrenan
los modelos de una vez y con tres semillas.

\textbf{¿El acierto del k-NN no está inflado?} El acierto dejando uno
fuera que muestra la aplicación sí puede estarlo (§12, nota al revisor).
En el artículo solo se dará el acierto en \textbf{testigos
independientes}.

\textbf{¿Qué pasa con la dependencia del color y de la cámara?} Se
documentan el equipo y la iluminación de cada imagen y se evalúa el
rendimiento entre fuentes distintas (entrenar con unas, probar con
otras). Los descriptores de tamaño van en unidades físicas (§11).

\textbf{¿Hay trabajos previos con deep learning en bioturbación?} Sí.
Hay que revisarlos y citarlos: el artículo de \emph{Applied Sciences}
11(16):7736 (2021) sobre deep learning en análisis icnológico y la
patente US12142024. La aportación propia es el \textbf{conjunto de datos
abierto}, la herramienta de anotación libre y la comparación con
concordancia humana.

\begin{center}\rule{0.5\linewidth}{0.5pt}\end{center}

\hypertarget{artuxedculo-3-validaciuxf3n-y-aplicaciuxf3n-paleoambiental}{%
\subsection{Artículo 3 --- Validación y aplicación
paleoambiental}\label{artuxedculo-3-validaciuxf3n-y-aplicaciuxf3n-paleoambiental}}

\hypertarget{quuxe9-describiruxe1-1}{%
\subsubsection{Qué describirá}\label{quuxe9-describiruxe1-1}}

\begin{enumerate}
\def\labelenumi{\arabic{enumi}.}
\tightlist
\item
  Cuánto difieren los índices automáticos de la estimación visual de
  varios icnólogos.
\item
  Registros continuos de bioturbación y de icnotaxones en testigos
  largos, correlacionados con indicadores geoquímicos y de temperatura
  (amplía el ejemplo de ODP 977 de 2020).
\end{enumerate}

\hypertarget{muxe9todos-que-usa-2}{%
\subsubsection{Métodos que usa}\label{muxe9todos-que-usa-2}}

\begin{longtable}[]{@{}ll@{}}
\toprule\noalign{}
Método & § \\
\midrule\noalign{}
\endhead
\bottomrule\noalign{}
\endlastfoot
Porcentaje, BI e ii automáticos & §3--§4 \\
BI e ii manuales en la aplicación & §4.3 \\
Porcentaje por icnotaxón e identificación asistida & §3, §12--§13 \\
Modelo de edad & §5 \\
Registro compuesto, registros externos y CSV interpolado & §15--§16 \\
\end{longtable}

\hypertarget{experimentos-necesarios-1}{%
\subsubsection{Experimentos
necesarios}\label{experimentos-necesarios-1}}

\begin{enumerate}
\def\labelenumi{\arabic{enumi}.}
\setcounter{enumi}{-1}
\tightlist
\item
  \textbf{Modelo de edad publicado:} el del sondeo de prueba es
  ilustrativo y no vale para interpretar.
\item
  \textbf{Experimento a ciegas:} de 3 a 5 icnólogos asignan BI e ii
  visuales a los mismos intervalos, sin ver el resultado automático. La
  aplicación ya permite registrarlo (columnas \texttt{BI\_manual} e
  \texttt{ii\_manual}).
\item
  \textbf{Estadística:}

  \begin{itemize}
  \tightlist
  \item
    Concordancia entre observadores (kappa de Cohen ponderado o
    coeficiente de correlación intraclase).
  \item
    Sesgo de cada observador frente al valor automático.
  \item
    Diagrama de Bland-Altman para el porcentaje.
  \end{itemize}
\item
  \textbf{Aplicación:}

  \begin{itemize}
  \tightlist
  \item
    Elegir testigos con buenas imágenes y datos de indicadores ya
    publicados (material de la datos privados).
  \item
    Correlación con desfase y prueba de significación que tenga en
    cuenta la autocorrelación de las series.
  \end{itemize}
\end{enumerate}

\hypertarget{preguntas-probables-de-revisores-1}{%
\subsubsection{Preguntas probables de
revisores}\label{preguntas-probables-de-revisores-1}}

\textbf{Si los observadores no coinciden entre sí, ¿con qué se valida el
método?} No se trata de que la máquina «acierte» a un experto. Se mide
cuánta variabilidad hay entre expertos y si el método automático cae
dentro de ella, con la ventaja de ser reproducible.

\textbf{¿Una medida 2D de la cara del testigo representa la
bioturbación?} Es la misma limitación de las estimaciones visuales, que
también se hacen sobre la cara expuesta. Se indica como limitación (§32)
y se mantiene la comparación en las mismas condiciones.

\textbf{¿Cómo afecta la incertidumbre del modelo de edad a las
correlaciones?} Se discute aparte. Si hace falta, se repite el análisis
con modelos de edad alternativos. El programa no propaga esa
incertidumbre (§5).

\begin{center}\rule{0.5\linewidth}{0.5pt}\end{center}

\hypertarget{artuxedculo-4-clasificaciuxf3n-automuxe1tica-de-intervalos-con-registros-multiproxy-descripciones-e-icnologuxeda}{%
\subsection{Artículo 4 --- Clasificación automática de intervalos con
registros multiproxy, descripciones e
icnología}\label{artuxedculo-4-clasificaciuxf3n-automuxe1tica-de-intervalos-con-registros-multiproxy-descripciones-e-icnologuxeda}}

\hypertarget{quuxe9-describiruxe1-2}{%
\subsubsection{Qué describirá}\label{quuxe9-describiruxe1-2}}

Un método reproducible para etiquetar intervalos (eventos como baja
oxigenación, aporte detrítico o corrientes de fondo, pero también
litologías o facies) en muchos sondeos a partir de registros numéricos
(geoquímica, foraminíferos, TOC, sedimentología\ldots) y de intervalos
descritos (litología, turbiditas\ldots), combinando el criterio experto
(intervalos y reglas) con aprendizaje supervisado, y su cartografía por
ventanas de edad. Incluye la icnología (bioturbación, BI, ii,
icnotaxones) como variables.

\hypertarget{muxe9todos-que-usa-3}{%
\subsubsection{Métodos que usa}\label{muxe9todos-que-usa-3}}

\begin{longtable}[]{@{}
  >{\raggedright\arraybackslash}p{(\columnwidth - 2\tabcolsep) * \real{0.7273}}
  >{\raggedright\arraybackslash}p{(\columnwidth - 2\tabcolsep) * \real{0.2727}}@{}}
\toprule\noalign{}
\begin{minipage}[b]{\linewidth}\raggedright
Método
\end{minipage} & \begin{minipage}[b]{\linewidth}\raggedright
§
\end{minipage} \\
\midrule\noalign{}
\endhead
\bottomrule\noalign{}
\endlastfoot
Importación, remuestreo, razones, log y escalado & §18.1--§18.2 \\
Capas de intervalos jerárquicas como variables (0/1) & §18.2--§18.3 \\
Etiquetas por intervalos y reglas & §18.3 \\
Clasificadores definidos por el usuario (familias objetivo, entradas,
modelo) & §18.6 \\
Bosque aleatorio, red neuronal (MLP), k-NN, Bayes ingenuo, k-medias &
§18.4, §18.5 \\
Clasificación manual en columnas y copia desde columnas automáticas &
§18.6 \\
Validación dejando un sondeo fuera & §18.8 \\
Mapas por edad e IDW & §19 \\
Bioturbación e índices como variables & §3--§4, §15, §18.1 \\
\end{longtable}

\hypertarget{quuxe9-se-puede-afirmar-hoy-2}{%
\subsubsection{Qué se puede afirmar
hoy}\label{quuxe9-se-puede-afirmar-hoy-2}}

\begin{itemize}
\tightlist
\item
  La implementación funciona y es reproducible (semilla) con datos
  sintéticos de verdad conocida (§18.12).
\item
  Primera prueba con datos reales (§21.8 bis): con registros públicos de
  ODP 977A, los tipos de ORL privados se reconocen mal (acierto
  equilibrado 36 \% por bloques). No se puede afirmar capacidad
  predictiva; faltan los datos geoquímicos propios (XRF) y más sondeos
  etiquetados.
\item
  Con los datos propios (V12, §21.8 ter): acierto equilibrado 17,7 \%
  dejando un sondeo fuera con 7 clases. La transferencia entre sondeos
  con variables y escáneres distintos es el problema principal que un
  artículo tendría que resolver (armonizar variables, más sondeos por
  tipo de variable).
\item
  V26 (misma datos privados): el clasificador creado desde los ajustes
  guardados da 18 \% de acierto equilibrado, coherente con la cifra
  anterior. En el proyecto sintético de entrenamiento, bosque aleatorio
  $\approx$ 91 \% y red neuronal $\approx$ 80 \% (verdad conocida; solo demuestra que la
  cadena funciona).
\end{itemize}

\hypertarget{experimentos-necesarios-2}{%
\subsubsection{Experimentos
necesarios}\label{experimentos-necesarios-2}}

\begin{enumerate}
\def\labelenumi{\arabic{enumi}.}
\tightlist
\item
  \textbf{Datos reales con interpretación independiente:} sondeos con
  eventos publicados (p.~ej. sapropeles datados, contornitas, niveles
  detríticos) que no hayan servido para definir las reglas.
\item
  \textbf{Etiquetas solo por intervalos} para entrenar y evaluar. Las
  reglas se evalúan aparte, como «modelo experto», frente a la misma
  verdad.
\item
  \textbf{Validación dejando sondeos enteros fuera} y, si hay
  suficientes, por regiones (bloques espaciales), para medir la
  transferencia a sondeos lejanos.
\item
  \textbf{Métricas:} acierto equilibrado, F1 por clase y matriz de
  confusión; incertidumbre por remuestreo de sondeos (bootstrap por
  sondeo).
\item
  \textbf{Comparaciones:} reglas expertas solas; bosque aleatorio; red
  neuronal; k-NN; Bayes; y el efecto de añadir la icnología (con y sin
  bioturbación/índices) y las capas descritas (con y sin litología) como
  variables. Vigilar que una capa descrita no sea casi la misma
  información que el objetivo (fuga de información), y declarar qué
  capas se usaron.
\item
  \textbf{Importancia por permutación} (por implementar) además de MDI.
\item
  \textbf{Sensibilidad} al paso de remuestreo, al escalado (global
  frente a por sondeo) y a la probabilidad mínima.
\item
  \textbf{Modelos de edad:} repetir los mapas con modelos alternativos o
  con la incertidumbre publicada de cada sondeo.
\item
  \textbf{Prueba final independiente:} reservar 2--3 sondeos que no se
  usan para elegir modelo ni ajustes, y dar en el artículo el acierto en
  ellos.
\end{enumerate}

\hypertarget{preguntas-probables-de-revisores-2}{%
\subsubsection{Preguntas probables de
revisores}\label{preguntas-probables-de-revisores-2}}

\textbf{¿El acierto no está inflado por la autocorrelación de las
series?} Lo limita: la validación deja sondeos enteros fuera o, si solo
hay uno, bloques contiguos con un margen (desde V0.37, el 5 \% del
intervalo etiquetado a cada lado del bloque de prueba se quita del
entrenamiento). Desde V0.37 el escalado global se recalcula en cada
pliegue solo con el entrenamiento y una familia objetivo nunca usa como
entrada a sus familias madre ni hijas (p.~ej. «sapropel» no usa «baja
oxigenación»). Con un único sondeo la autocorrelación más allá del
margen sigue pudiendo inflar el acierto; con el escalado por sondeo,
cada sondeo de prueba se tipifica con sus propias medias (sin usar
etiquetas). Lo fiable es el acierto en sondeos reservados (§18.8).

\textbf{Si las etiquetas salen de reglas, ¿el modelo no aprende solo las
reglas?} Sí, y así se declara (§18.12). Para la evaluación se usan
intervalos con interpretación independiente; las reglas solo sirven para
proponer etiquetas o como referencia experta.

\textbf{¿Por qué no deep learning?} El programa incluye una red neuronal
(perceptrón multicapa con parada temprana, §18.5) y se compara con el
bosque aleatorio, k-NN y Bayes en el mismo esquema de validación. Con
pocas decenas o cientos de intervalos etiquetados, los árboles suelen
ser más robustos y dan una medida de importancia; redes más profundas
(recurrentes, convolucionales sobre la serie) necesitarían muchos más
sondeos etiquetados de los disponibles.

\textbf{Probasteis varios modelos: ¿no estáis eligiendo el que mejor
sale por azar?} Se informa de todos los modelos probados con la misma
validación y el acierto final se da en sondeos reservados que no
intervinieron en la elección (experimento 9).

\textbf{¿La interpolación del mapa es fiable?} La interpolación es
visual y exploratoria (§19, §19.5): ningún método estima la
incertidumbre de forma útil con pocos sondeos (el kriging da una
varianza, pero su variograma con \textasciitilde10 sondeos es muy
incierto) y ninguno tiene en cuenta barreras geográficas. Las
conclusiones se basan en los valores de cada sondeo, que también se
exportan. Los fondos de los mapas (Esri, GEBCO; §19.2) son solo contexto
visual: no entran en ningún cálculo, y su fuente debe citarse en el pie
de la figura (el programa la escribe en el mapa).

\textbf{¿Cómo entran datos categóricos como la litología en un modelo
numérico?} Como columnas 0/1 por categoría, sin escalar (§18.2). La
ausencia de descripción vale 0, igual que «no presente»; se declara como
limitación y conviene usar solo sondeos con la capa descrita entera.

\textbf{¿Cómo se combinan proxies medidos con técnicas o calibraciones
distintas?} Con el z-score por sondeo (§18.2) se eliminan diferencias de
nivel; se documenta qué escalado se usó y se muestra su efecto
(experimento 7).

\begin{center}\rule{0.5\linewidth}{0.5pt}\end{center}

\hypertarget{artuxedculo-5-nota-breve-en-espauxf1ol-geogaceta}{%
\subsection{Artículo 5 --- Nota breve en español
(Geogaceta)}\label{artuxedculo-5-nota-breve-en-espauxf1ol-geogaceta}}

\hypertarget{quuxe9-describiruxe1-3}{%
\subsubsection{Qué describirá}\label{quuxe9-describiruxe1-3}}

Presentación de SedCoreTool a la comunidad geológica española: qué hace,
para quién es y un ejemplo con un testigo de datos públicos (IODP
U1385I-1H, margen ibérico) y el proyecto de entrenamiento. Sin
resultados inéditos.

\hypertarget{quuxe9-se-puede-afirmar-hoy-3}{%
\subsubsection{Qué se puede afirmar
hoy}\label{quuxe9-se-puede-afirmar-hoy-3}}

Lo mismo que en el artículo 1 (funcionalidad y reproducibilidad), sin
cifras de validación nuevas.

\hypertarget{quuxe9-falta}{%
\subsubsection{Qué falta}\label{quuxe9-falta}}

\begin{itemize}
\tightlist
\item[$\square$]
  Release con DOI y aplicación accesible en línea (GitHub Pages).
\item[$\square$]
  Figura: imagen tratada + columna compuesta + mapa de ejemplo.
\item[$\square$]
  Comprobar formato y extensión en las normas de la revista.
\end{itemize}

\hypertarget{preguntas-probables-de-revisores-3}{%
\subsubsection{Preguntas probables de
revisores}\label{preguntas-probables-de-revisores-3}}

\textbf{¿Qué aporta frente a hacerlo a mano en Photoshop/ImageJ y una
hoja de cálculo?} Un flujo único, reproducible y gratuito desde la
imagen hasta series, mapas y figuras, con el proyecto guardado para
repetir el análisis.

\begin{center}\rule{0.5\linewidth}{0.5pt}\end{center}

\hypertarget{lista-de-comprobaciuxf3n-antes-de-enviar-cualquier-artuxedculo}{%
\subsection{Lista de comprobación antes de enviar cualquier
artículo}\label{lista-de-comprobaciuxf3n-antes-de-enviar-cualquier-artuxedculo}}

\begin{itemize}
\tightlist
\item[$\square$]
  Cada método citado en el texto tiene su sección en
  \texttt{TECHNICAL\_DOCUMENTATION.md}, y coincide con el código de la
  versión citada.
\item[$\square$]
  Versión del software citada con DOI de Zenodo.
\item[$\square$]
  Parámetros usados guardados (archivo de proyecto) y publicados como
  material suplementario.
\item[$\square$]
  Datos y máscaras depositados (Zenodo o PANGAEA) con licencia.
\item[$\square$]
  Limitaciones declaradas (§32).
\item[$\square$]
  Referencias comprobadas contra los originales.
\item[$\square$]
  Coautorías y agradecimientos acordados.
\end{itemize}

\newpage

\hypertarget{borrador-1-software-joss-en-ingluxe9s}{%
\section{Borrador 1 --- Software (JOSS, en
inglés)}\label{borrador-1-software-joss-en-ingluxe9s}}

\hypertarget{summary}{%
\subsection{Summary}\label{summary}}

Bioturbation intensity is a key parameter of the ichnofabric approach
and is routinely estimated in sediment cores with semi-quantitative
visual schemes such as the Bioturbation Index (Reineck 1963; Taylor y
Goldring 1993) and the ichnofabric index (Droser y Bottjer 1986). Visual
estimation is subjective, slow, and usually discontinuous along the
core. IchAnTool converts segmented core images into continuous, depth-
or age-referenced series of bioturbated-surface percentage, per-class
coverage (e.g.~per ichnotaxon), and the corresponding semi-quantitative
indices, ready to be correlated with geochemical, mineralogical, or
logging data.

\hypertarget{statement-of-need}{%
\subsection{Statement of need}\label{statement-of-need}}

Digital image treatment has improved the visibility and quantification
of trace fossils in cores (Dorador y Rodríguez-Tovar 2014, 2018), but
turning treated images into continuous data series remained a manual,
repetitive task. A first version (IchAnTool), a single script, was
introduced in Casanova-Arenillas, Rodríguez-Tovar, y Martínez-Ruiz
(2020). SedCoreTool is a complete rewrite: an installable, tested Python
package with a command-line interface and API, and a browser application
that needs no installation, aimed at sedimentologists, ichnologists,
palaeoceanographers, and core-analysis practitioners in academia and
industry. {[}TODO: one paragraph on existing alternatives and why
IchAnTool is needed.{]}

\hypertarget{functionality}{%
\subsection{Functionality}\label{functionality}}

\begin{itemize}
\tightlist
\item
  Input: treated images (PNG, TIFF, JPEG) with bioturbation painted in a
  single colour, colour-coded classes, or integer label masks.
\item
  Row-wise bioturbated percentage, aggregated into user-defined depth
  bins.
\item
  Per-class percentages for colour-coded ichnotaxa.
\item
  Configurable index schemes (BI and ii provided).
\item
  Depth-to-age conversion with user age models.
\item
  CSV/Excel export, quick-look plots, and a modular segmentation
  interface for future learned models.
\item
  Results reproduce those of the published version exactly (regression
  test).
\item
  Browser application: projects with several cores; numeric records of
  any kind (XRF, foraminifera, TOC, colour) in depth or age; families of
  described intervals classified by hand; user-defined classifiers
  (random forest, multilayer perceptron, k-nearest neighbours, naive
  Bayes, k-means) validated leaving whole cores out; maps by age window
  with several projections and interpolation methods; Grapher-like
  plots; publication figures; a synthetic training project with a guided
  tour.
\end{itemize}

\hypertarget{implementation}{%
\subsection{Implementation}\label{implementation}}

Both implementations share the same definitions. A label mask (0 =
matrix, 1..n = objects or classes) is restricted to an analysis area.
The bioturbated percentage of each depth interval of \emph{N} pixel rows
is the ratio of bioturbated to valid pixels. Excluded pixels (cracks,
voids, samples; currently in the web application only) leave both terms.
Each interval is assigned the depth of its centre, interpolated linearly
between the top and bottom depths of the analysis area. Percentages map
to BI and ii through half-open tiers, and ages come from linear
interpolation of a user age model.

The web application adds four steps before the mask is final:

\begin{itemize}
\tightlist
\item
  image enhancement: levels, gamma, contrast, saturation, histogram
  equalisation and unsharp-mask local contrast;
\item
  candidate detection: adaptive mean-C thresholding, 3 × 3 morphological
  opening/closing and 8-connected components with a minimum size;
\item
  interactive editing: colour region growing, brush, lasso and an
  exclusion layer;
\item
  assisted ichnotaxon identification: a distance-weighted
  \emph{k}-nearest-neighbour classifier (\emph{k} = 5) on twelve
  standardised shape, size (physical units) and colour descriptors, with
  a confidence threshold, a novelty gate based on the leave-one-out
  nearest-neighbour distance and exportable models. Optionally, the
  descriptors can be replaced by the embedding of a frozen pre-trained
  vision model (DINOv2; the classifier still learns only from the user's
  labels);
\item
  learned segmentation (experimental): export of images and expert
  masks, a U-Net trained with the Python package or from the desktop
  program on the user's GPU, and in-app inference with ONNX Runtime
  (probability threshold and doubtful zones). No trained model is
  distributed yet.
\end{itemize}

All computations except optional U-Net training are deterministic.
Project files store the original image and every parameter. A full
technical reference, with formulas and default values, is distributed
with the code (\texttt{docs/}).

{[}TODO: figure --- example core, mask and output curves
(\texttt{examples/}).{]}

\hypertarget{acknowledgements}{%
\subsection{Acknowledgements}\label{acknowledgements}}

{[}TODO{]}

\hypertarget{references}{%
\subsection{References}\label{references}}

\newpage

\hypertarget{borrador-2-dataset-y-segmentaciuxf3n}{%
\section{Borrador 2 --- Dataset y
segmentación}\label{borrador-2-dataset-y-segmentaciuxf3n}}

\hypertarget{paper-2-open-annotated-dataset-and-automatic-segmentation-of-bioturbation-in-core-images}{%
\subsection{Paper 2 --- Open annotated dataset and automatic
segmentation of bioturbation in core
images}\label{paper-2-open-annotated-dataset-and-automatic-segmentation-of-bioturbation-in-core-images}}

\textbf{Target:} Earth System Science Data (data focus) or Computers \&
Geosciences / Artificial Intelligence in Geosciences (method focus).
Decide once the dataset size and model results are known.

\textbf{Plain-language guide (Spanish):}
\texttt{publications/GUIA\_APRENDIZAJE\_AUTOMATICO.md}.

\textbf{Status (V0.17, 05-11-2026; training plan in §3.5; training
inside the desktop program and vision-model descriptors for ichnotaxa in
§3.6):} the full workflow is implemented in SedCoreTool and tested on
synthetic data (export $\rightarrow$ tiles split by core $\rightarrow$ training $\rightarrow$ evaluation on
whole images of independent cores $\rightarrow$ use in the app). \textbf{No model
has been trained on real cores yet; every number below is a
placeholder.} Methods text describes the code exactly; keep it in sync
with \texttt{docs/TECHNICAL\_DOCUMENTATION.md} §9.

\hypertarget{working-title}{%
\subsubsection{Working title}\label{working-title}}

An open, annotated dataset of sediment-core images and a U-Net baseline
for automatic segmentation of bioturbation {[}TODO: ``and ichnotaxa'' if
multi-class results are good enough{]}

\hypertarget{provisional-abstract}{%
\subsubsection{Provisional abstract}\label{provisional-abstract}}

Quantifying bioturbation in cores still depends on manual selection of
trace fossils on images, the main bottleneck of image-based ichnofabric
analysis. We release {[}N{]} images from {[}M{]} cores
({[}sites/expeditions{]}) with expert masks of bioturbation {[}and
ichnotaxa{]}, annotated in the open program SedCoreTool, together with
an annotation protocol and inter-annotator agreement on a subset. We
train a U-Net on cores split into training, validation and test sets, so
that all scores are measured on cores never seen during training, and
compare it with classical thresholding workflows used in ichnology. On
the test cores the U-Net reaches {[}Dice{]} (pixel) and an error of
{[}x{]} percentage points in bioturbation per {[}b{]}-cm interval, with
BI agreement {[}y{]} \%, against {[}Dice\_base{]} and {[}x\_base{]} for
the best classical workflow and {[}Dice\_human{]} between experts. Data,
code and the trained model are openly released; the model runs inside
SedCoreTool without a GPU.

\hypertarget{introduction}{%
\subsubsection{1. Introduction}\label{introduction}}

\begin{itemize}
\tightlist
\item
  Image-based ichnofabric analysis: enhancement + manual selection
  (Dorador et al., 2014; Miguez-Salas et al., 2019); percentage per
  interval and BI/ii (Casanova-Arenillas et al., 2020).
\item
  Bottleneck: manual selection is slow and subjective; thresholding
  depends on contrast between infill and host sediment and fails with
  low contrast, mottling, cracks or light infills.
\item
  Deep learning for segmentation (Ronneberger et al., 2015) has been
  applied to bioturbation {[}TODO: review --- Applied Sciences
  11(16):7736 (2021); Kikuchi and Naruse (2024), Dice 0.892 for
  \emph{Phycosiphon} --- verify authors, titles and journals before
  citing; patent US12142024: mention only if relevant to related
  work{]}.
\item
  Gap: no open, multi-site annotated dataset; scores often reported on
  random tiles, inflated by spatial autocorrelation; trained models
  rarely usable by ichnologists.
\item
  Contributions: (1) open dataset + protocol + agreement; (2) U-Net
  baseline evaluated on independent cores against classical workflows
  and human agreement; (3) integration in an open program (annotation,
  training data export, inference, correction).
\end{itemize}

\hypertarget{data}{%
\subsubsection{2. Data}\label{data}}

\begin{itemize}
\tightlist
\item
  Sources {[}TODO{]}: IODP/ODP core images (public), {[}UGR collections
  --- permission and licence{]}.
\item
  Per image: core, section, top and bottom depth, resolution (px/cm),
  camera/scanner and lighting {[}TODO: add equipment to the manifest or
  a side table{]}.
\item
  Annotation protocol {[}TODO: write
  \texttt{publications/02\_dataset\_segmentation/protocol.md}{]}: what
  counts as a trace, artefacts (cracks, voids, sampling holes, drilling
  disturbance) excluded with the exclusion layer, ichnotaxa criteria,
  minimum size.
\item
  Annotation tool: SedCoreTool (wand, brush, lasso, polygon, exclusion;
  candidates from classical detection reviewed one by one). Only traces
  with a class confirmed by the annotator are used.
\item
  Inter-annotator agreement: $\geq$ 2 experts on {[}n{]} images; IoU, Dice
  and difference in \% per interval.
\end{itemize}

\hypertarget{exported-format-sedcoretool-results-u-net-training-data}{%
\paragraph{\texorpdfstring{Exported format (SedCoreTool \emph{Results $\rightarrow$
U-Net training
data})}{Exported format (SedCoreTool Results $\rightarrow$ U-Net training data)}}\label{exported-format-sedcoretool-results-u-net-training-data}}

One ZIP per project:
\texttt{images/\textless{}core\textgreater{}\_\_\textless{}section\textgreater{}.png}
(original analysis area, no enhancement),
\texttt{masks/\textless{}core\textgreater{}\_\_\textless{}section\textgreater{}.png}
(8-bit: 0 = matrix or reviewed artefact, k = class k, 255 = ignore:
excluded zones and traces without a confirmed class) and
\texttt{manifest.json} (classes; per image core, section, top, bottom,
units, px/cm, size and counts of reviewed and unreviewed objects).
Unreviewed traces are ignored rather than counted as matrix, because an
unconfirmed trace is not a reliable negative.

\hypertarget{methods}{%
\subsubsection{3. Methods}\label{methods}}

\hypertarget{data-preparation-sedcoretool-unet-prepare}{%
\paragraph{\texorpdfstring{3.1 Data preparation
(\texttt{sedcoretool\ unet\ prepare})}{3.1 Data preparation (sedcoretool unet prepare)}}\label{data-preparation-sedcoretool-unet-prepare}}

\begin{itemize}
\tightlist
\item
  All images resampled to a common resolution ρ (default: median px/cm
  of the dataset); images bilinear, masks nearest neighbour.
\item
  Tiles of T = 512 px with 64-px overlap; tiles with \textless{} 50 \%
  valid (non-ignored) pixels discarded; areas outside the image marked
  as ignore.
\item
  \textbf{Split by core}: all tiles of a core in the same set. Test
  cores chosen {[}by site / at random with seed 42{]}; never used for
  training, early stopping or threshold choice.
\item
  Binary target (bioturbation vs matrix) {[}and multi-class variant{]}.
\end{itemize}

\hypertarget{model-and-training-sedcoretool-unet-train}{%
\paragraph{\texorpdfstring{3.2 Model and training
(\texttt{sedcoretool\ unet\ train})}{3.2 Model and training (sedcoretool unet train)}}\label{model-and-training-sedcoretool-unet-train}}

\begin{itemize}
\tightlist
\item
  U-Net with ResNet-34 encoder pre-trained on ImageNet
  (segmentation-models-pytorch).
\item
  Loss: binary cross-entropy with logits + soft Dice, ignoring 255
  pixels (multi-class: cross-entropy

  \begin{itemize}
  \tightlist
  \item
    mean Dice of classes $\geq$ 1).
  \end{itemize}
\item
  Augmentation: horizontal and vertical flips; contrast ×U(0.85, 1.15),
  brightness +U(−0.08, 0.08), gamma U(0.8, 1.25). No 90° rotations:
  trace orientation relative to bedding is diagnostic.
\item
  AdamW (lr 3·10⁻⁴, weight decay 10⁻⁴), cosine schedule, batch 8, $\leq$ 60
  epochs, early stopping (patience 12) on validation Dice. Seed 42.
  Hardware {[}TODO: GPU model, training time{]}.
\end{itemize}

\hypertarget{inference-same-rules-in-python-and-in-the-app}{%
\paragraph{3.3 Inference (same rules in Python and in the
app)}\label{inference-same-rules-in-python-and-in-the-app}}

\begin{itemize}
\tightlist
\item
  Whole analysis area resampled to ρ, cut into overlapping tiles,
  normalised with ImageNet mean and SD, sigmoid of the logits, tiles
  merged with a linear ramp weight over the overlap, resampled back.
  Python and JavaScript implementations are checked against the same
  reference values.
\item
  Candidates: P $\geq$ 0.5 (threshold fixed on validation cores, not on test
  cores), then the same morphological cleaning and minimum size as the
  classical workflow.
\item
  In the program, users can show ``doubtful'' pixels (\textbar P −
  0.5\textbar{} \textless{} 0.2) and correct the result with the editing
  tools; corrected images can be exported again to enlarge the dataset.
\end{itemize}

\hypertarget{evaluation-sedcoretool-unet-evaluate}{%
\paragraph{\texorpdfstring{3.4 Evaluation
(\texttt{sedcoretool\ unet\ evaluate})}{3.4 Evaluation (sedcoretool unet evaluate)}}\label{evaluation-sedcoretool-unet-evaluate}}

\begin{itemize}
\tightlist
\item
  On \textbf{whole images} of the test cores, not tiles.
\item
  Pixel: IoU, Dice, precision, recall (ignored pixels excluded); pooled
  over images (TP, FP, FN summed).
\item
  Per interval of b cm {[}1 cm and \ldots{]}: absolute error and bias of
  bioturbation \%, agreement of BI (Taylor and Goldring, 1993) and ii
  (Droser and Bottjer, 1986) classes.
\item
  Baselines on the same test images: adaptive threshold (SedCoreTool
  default); CLAHE + threshold or K-means + hole filling (Miguez-Salas et
  al., 2019); Dorador et al.~(2014) enhancement + threshold. {[}TODO:
  implement the classical detections in the Python package, or export
  their masks from the app, so that baselines are scored with the same
  code.{]}
\item
  Human ceiling: agreement between annotators.
\item
  {[}Optional{]} Cross-source test: train on one source (e.g.~IODP
  scans), test on another.
\end{itemize}

\hypertarget{training-protocol-and-experiments-plan-v0.16}{%
\paragraph{3.5 Training protocol and experiments (plan,
V0.16)}\label{training-protocol-and-experiments-plan-v0.16}}

Run in this order; each step uses only commands that exist in
SedCoreTool V0.15.

\begin{enumerate}
\def\labelenumi{\arabic{enumi}.}
\tightlist
\item
  \textbf{Pilot (3--5 cores).} Annotate, export, prepare, train for a
  few epochs and evaluate, to check the whole chain on real images and
  measure annotation time per metre of core. Pilot cores can be re-used
  later only if their annotation follows the final protocol.
\item
  \textbf{Dataset (target 15--20 cores, $\geq$ 2 sources/sites).} Annotate
  following the written protocol. Record source, equipment, resolution
  and lithology per core.
\item
  \textbf{Lock the test set before training.} Choose the test cores ($\approx$
  20 \%, covering every source and lithology) and write their names in
  \texttt{dataset.json} (\texttt{-\/-test}). They are not looked at
  again until the final evaluation: no training, no early stopping, no
  threshold or parameter choice on them.
\item
  \textbf{Validation cores ($\approx$ 20 \%)} for early stopping and for
  choosing the probability threshold: evaluate the validation cores
  (\texttt{evaluate\ -\/-cores\ …}) at thresholds 0.3--0.7 and keep the
  one with the best Dice; that value is then fixed for the test set.
\item
  \textbf{Main model:} ResNet-34 pre-trained on ImageNet, settings of
  §3.2, seed 42. Repeat with 3 seeds and report mean $\pm$ SD to show the
  variability due to training.
\item
  \textbf{Baselines and human ceiling} on the same test images with the
  same scoring code (§3.4).
\item
  \textbf{Learning curve:} retrain with 25, 50, 75 and 100 \% of the
  training cores (same test set) to show how many annotated cores are
  needed.
\item
  \textbf{Cross-source test (if $\geq$ 2 sources):} train on one source, test
  on the other.
\item
  \textbf{Correction loop (optional, reported separately):} the model's
  doubtful zones guide the expert's corrections on new cores; corrected
  cores are added to training. Annotation time with and without the
  model's proposal is measured. Cores annotated this way never go to the
  test set, because the proposal could bias the expert.
\item
  \textbf{Release:} model \texttt{.onnx} + \texttt{.json} (with split,
  seed and history), dataset and protocol with DOIs.
\end{enumerate}

\hypertarget{software-used-for-the-experiments-v0.17}{%
\paragraph{3.6 Software used for the experiments
(V0.17)}\label{software-used-for-the-experiments-v0.17}}

\begin{itemize}
\tightlist
\item
  \textbf{Training inside SedCoreTool (desktop).} The same Python code
  as the command line, run from the program on the annotator's computer
  (NVIDIA CUDA, Apple MPS or CPU). The user marks each core as
  validation, test or automatic; the program exports the training data,
  prepares tiles, trains, evaluates on the test cores and loads the
  model. Training can be stopped and resumed from a per-epoch checkpoint
  (models reported in the paper are trained without interruption). This
  matters for the dataset paper: annotators without programming
  experience can reproduce the training and the correction loop (§3.5,
  step 9).
\item
  \textbf{Ichnotaxon identification (secondary aim).} Two inputs for the
  same classifiers (distance-weighted k-NN with novelty rejection,
  Random Forest, naive Bayes): (a) 12 shape and colour descriptors per
  trace;

  \begin{enumerate}
  \def\labelenumi{(\alph{enumi})}
  \setcounter{enumi}{1}
  \tightlist
  \item
    the 384-value embedding of a frozen, self-supervised vision model
    (DINOv2 ViT-S/14; Oquab et al., 2024) of a square window around each
    trace (25 \% context per side, original image, 224 × 224 px, unit
    length). The vision model is not fine-tuned; only the classifier
    learns from the annotations. Both inputs are compared on the same
    independent test cores (balanced accuracy, per-class F1, confusion
    matrix) and against agreement between experts. {[}TODO: decide
    whether identification goes in this paper or in paper 3.{]}
  \end{enumerate}
\end{itemize}

\hypertarget{results-todo-no-results-yet}{%
\subsubsection{4. Results {[}TODO --- no results
yet{]}}\label{results-todo-no-results-yet}}

\begin{longtable}[]{@{}
  >{\raggedright\arraybackslash}p{(\columnwidth - 14\tabcolsep) * \real{0.1250}}
  >{\raggedright\arraybackslash}p{(\columnwidth - 14\tabcolsep) * \real{0.1250}}
  >{\raggedright\arraybackslash}p{(\columnwidth - 14\tabcolsep) * \real{0.1250}}
  >{\raggedright\arraybackslash}p{(\columnwidth - 14\tabcolsep) * \real{0.1250}}
  >{\raggedright\arraybackslash}p{(\columnwidth - 14\tabcolsep) * \real{0.1250}}
  >{\raggedright\arraybackslash}p{(\columnwidth - 14\tabcolsep) * \real{0.1250}}
  >{\raggedright\arraybackslash}p{(\columnwidth - 14\tabcolsep) * \real{0.1250}}
  >{\raggedright\arraybackslash}p{(\columnwidth - 14\tabcolsep) * \real{0.1250}}@{}}
\toprule\noalign{}
\begin{minipage}[b]{\linewidth}\raggedright
Method
\end{minipage} & \begin{minipage}[b]{\linewidth}\raggedright
IoU
\end{minipage} & \begin{minipage}[b]{\linewidth}\raggedright
Dice
\end{minipage} & \begin{minipage}[b]{\linewidth}\raggedright
Precision
\end{minipage} & \begin{minipage}[b]{\linewidth}\raggedright
Recall
\end{minipage} & \begin{minipage}[b]{\linewidth}\raggedright
MAE \% (b cm)
\end{minipage} & \begin{minipage}[b]{\linewidth}\raggedright
BI agreement
\end{minipage} & \begin{minipage}[b]{\linewidth}\raggedright
ii agreement
\end{minipage} \\
\midrule\noalign{}
\endhead
\bottomrule\noalign{}
\endlastfoot
Adaptive threshold & & & & & & & \\
CLAHE + threshold / K-means (Miguez-Salas et al., 2019) & & & & & & & \\
Dorador et al.~(2014) + threshold & & & & & & & \\
U-Net (this work) & & & & & & & \\
Second annotator & & & & & & & \\
\end{longtable}

Identification (test cores):

\begin{longtable}[]{@{}lllll@{}}
\toprule\noalign{}
Input & Classifier & Balanced accuracy & Macro F1 & Rejected (novel) \\
\midrule\noalign{}
\endhead
\bottomrule\noalign{}
\endlastfoot
12 shape/colour descriptors & k-NN & & & \\
DINOv2 embedding & k-NN & & & \\
best of RF / Bayes & & & & \\
Second expert & --- & & & --- \\
\end{longtable}

Figures: (1) dataset map and examples; (2) workflow; (3) examples of
predictions, errors and doubtful zones on test cores; (4) bioturbation
\% logs, expert vs U-Net vs baseline; (5) error by contrast class /
lithology.

\hypertarget{discussion-todo}{%
\subsubsection{5. Discussion {[}TODO{]}}\label{discussion-todo}}

\begin{itemize}
\tightlist
\item
  Where the U-Net fails (light infills, mottled background, cracks) and
  how users correct it.
\item
  Effect of resolution and source; how many annotated cores are needed
  (learning curve).
\item
  Limits: 2-D surfaces; annotation subjectivity; class imbalance for
  rare ichnotaxa.
\end{itemize}

\hypertarget{code-and-data-availability}{%
\subsubsection{6. Code and data
availability}\label{code-and-data-availability}}

\begin{itemize}
\tightlist
\item
  Code: SedCoreTool (GPL-3.0) {[}repository URL and Zenodo DOI of the
  version used{]}.
\item
  Dataset: {[}Zenodo/PANGAEA DOI{]}, licence {[}CC BY 4.0?{]}, including
  masks, manifest and protocol.
\item
  Trained model: \texttt{.onnx} + \texttt{.json} {[}DOI{]}; runs in
  SedCoreTool (desktop, browser) and with
  \texttt{sedcoretool\ unet\ predict}.
\end{itemize}

\hypertarget{needs}{%
\subsubsection{Needs}\label{needs}}

\begin{itemize}
\tightlist
\item[$\boxtimes$]
  Annotation, export, training, evaluation and in-app inference
  implemented (V0.15).
\item[$\square$]
  Permission and licence for the images (IODP images are public; UGR
  treated images: ask).
\item[$\square$]
  Annotation protocol written.
\item[$\square$]
  Annotation by $\geq$ 2 experts on a subset to measure agreement.
\item[$\square$]
  Classical baselines scored with the same evaluation code.
\item[$\square$]
  Equipment and lighting per image recorded.
\item[$\square$]
  Pilot with 3--5 cores (§3.5, step 1).
\item[$\square$]
  Test cores chosen and locked before any training.
\item[$\square$]
  Training on a GPU (a single consumer GPU is enough), 3 seeds, learning
  curve.
\item[$\square$]
  Compare shape descriptors vs DINOv2 embeddings for ichnotaxa on test
  cores.
\item[$\square$]
  Verify all references in §1.
\end{itemize}

\hypertarget{references-to-complete-and-verify}{%
\subsubsection{References (to complete and
verify)}\label{references-to-complete-and-verify}}

\begin{itemize}
\tightlist
\item
  Casanova-Arenillas, S. et al.~(2020). \emph{Computers \& Geosciences}
  136, 104407. doi:10.1016/j.cageo.2020.104407
\item
  Dorador, J., Rodríguez-Tovar, F. J. and IODP Expedition 339 Scientists
  (2014). \emph{Facies} 60.
\item
  Droser, M. L. and Bottjer, D. J. (1986). \emph{Journal of Sedimentary
  Petrology} 56.
\item
  He, K., Zhang, X., Ren, S. and Sun, J. (2016). Deep residual learning
  for image recognition. CVPR.
\item
  Miguez-Salas, O., Dorador, J. and Rodríguez-Tovar, F. J. (2019).
  \emph{Marine Geology} 413.
\item
  Oquab, M. et al.~(2024). DINOv2: learning robust visual features
  without supervision. \emph{Transactions on Machine Learning Research}.
  {[}verify{]}
\item
  Ronneberger, O., Fischer, P. and Brox, T. (2015). U-Net: convolutional
  networks for biomedical image segmentation. MICCAI, LNCS 9351.
\item
  Taylor, A. M. and Goldring, R. (1993). \emph{Journal of the Geological
  Society} 150.
\item
  {[}TODO{]} Applied Sciences 11(16):7736 (2021); Kikuchi and Naruse
  (2024).
\end{itemize}

\newpage

\hypertarget{borrador-3-validaciuxf3n-y-aplicaciuxf3n}{%
\section{Borrador 3 --- Validación y
aplicación}\label{borrador-3-validaciuxf3n-y-aplicaciuxf3n}}

\hypertarget{paper-3-validation-and-palaeoenvironmental-application}{%
\subsection{Paper 3 --- Validation and palaeoenvironmental
application}\label{paper-3-validation-and-palaeoenvironmental-application}}

\textbf{Target:} Palaeogeography, Palaeoclimatology, Palaeoecology.
Alternatives: Ichnos, Marine Geology, Sedimentary Geology.

\hypertarget{working-title-1}{%
\subsubsection{Working title}\label{working-title-1}}

Continuous, objective bioturbation records from core images: validation
against visual indices and application to {[}region/period{]} {[}TODO{]}

\hypertarget{idea}{%
\subsubsection{Idea}\label{idea}}

\begin{itemize}
\tightlist
\item
  Validation: compare IchAnTool percentages with blind visual BI/ii
  estimates by several ichnologists on the same intervals; quantify
  subjectivity and bias of visual schemes.
\item
  Application: continuous bioturbation and per-ichnotaxon records along
  long cores, correlated with geochemical/redox and SST proxies
  (e.g.~Western Mediterranean during the Last Glacial Cycle; extends the
  ODP 977 example of Casanova-Arenillas et al., 2020).
\end{itemize}

\hypertarget{outline}{%
\subsubsection{Outline}\label{outline}}

\begin{enumerate}
\def\labelenumi{\arabic{enumi}.}
\tightlist
\item
  Introduction
\item
  Material --- cores, images, age models, proxy data
\item
  Methods --- IchAnTool workflow, observer experiment, statistics
\item
  Results --- agreement and bias; continuous records
\item
  Discussion --- what continuous records reveal that discrete estimates
  miss
\item
  Conclusions
\end{enumerate}

\hypertarget{needs-1}{%
\subsubsection{Needs}\label{needs-1}}

\begin{itemize}
\tightlist
\item[$\square$]
  Choose cores with good images and existing proxy data (own material).
\item[$\square$]
  Recruit 3--5 observers for the blind test.
\end{itemize}

\newpage

\hypertarget{borrador-4-clasificaciuxf3n-multiproxy}{%
\section{Borrador 4 --- Clasificación
multiproxy}\label{borrador-4-clasificaciuxf3n-multiproxy}}

\hypertarget{paper-4-automatic-classification-of-core-intervals-from-multiproxy-records}{%
\subsection{Paper 4 --- Automatic classification of core intervals from
multiproxy
records}\label{paper-4-automatic-classification-of-core-intervals-from-multiproxy-records}}

\textbf{Target:} Computers \& Geosciences. Alternatives: Artificial
Intelligence in Geosciences, Applied Computing and Geosciences; Marine
Geology if the emphasis is the palaeoceanographic result.

\hypertarget{working-title-2}{%
\subsubsection{Working title}\label{working-title-2}}

Reproducible classification of core intervals (sapropels, organic-rich
layers, turbidites\ldots) from multiproxy records with user-defined
machine-learning classifiers {[}TODO{]}

\hypertarget{provisional-abstract-1}{%
\subsubsection{Provisional abstract}\label{provisional-abstract-1}}

Interval interpretation in cores (e.g.~low-oxygen events, detrital
pulses, climate stages) is done by experts core by core and is hard to
reproduce across many cores and laboratories. We present a workflow in
SedCoreTool in which experts define families of intervals and their
classes, label some cores, and train classifiers (random forest,
multilayer perceptron, k-nearest neighbours, naive Bayes) on resampled
multiproxy records (XRF, foraminifera, TOC, colour, bioturbation).
Models are validated leaving whole cores out and tested on {[}N{]} cores
withheld from model selection. {[}TODO: results --- balanced accuracy
per family, per class F1, transfer between basins, effect of adding
ichnology.{]}

\hypertarget{outline-1}{%
\subsubsection{Outline}\label{outline-1}}

\begin{enumerate}
\def\labelenumi{\arabic{enumi}.}
\tightlist
\item
  Introduction --- expert interval interpretation; reproducibility
  across cores; prior ML in core logging {[}TODO review{]}.
\item
  Data --- cores, proxies, age models, described intervals (independent
  of the rules), harmonisation of variables.
\item
  Methods --- resampling and scaling (§21.2); families and labels
  (§21.3); classifiers (§21.4--§21.4 ter); validation leaving one core
  out (§21.5); held-out cores for the final test; permutation importance
  {[}TODO implement{]}.
\item
  Results --- comparison of models and inputs; maps by age window
  (§21.6) as illustration only.
\item
  Discussion --- transfer between cores with different scanners; leakage
  between families; limits.
\item
  Software and data availability --- SedCoreTool release (DOI), project
  files as supplementary material.
\end{enumerate}

\hypertarget{needs-2}{%
\subsubsection{Needs}\label{needs-2}}

\begin{itemize}
\tightlist
\item[$\square$]
  Cores with independently described events (published sapropels/ORL,
  turbidites) not used to define rules.
\item[$\square$]
  Harmonised variable set across cores (same XRF ratios or calibrated
  values).
\item[$\square$]
  Held-out cores option in the interface (to avoid model-selection
  bias).
\item[$\square$]
  Permutation importance (also valid for the neural network).
\end{itemize}

\newpage

\hypertarget{borrador-5-nota-en-geogaceta}{%
\section{Borrador 5 --- Nota en
Geogaceta}\label{borrador-5-nota-en-geogaceta}}

\hypertarget{artuxedculo-5-nota-breve-sedcoretool-una-herramienta-abierta-para-el-anuxe1lisis-de-sondeos}{%
\subsection{Artículo 5 --- Nota breve: SedCoreTool, una herramienta
abierta para el análisis de
sondeos}\label{artuxedculo-5-nota-breve-sedcoretool-una-herramienta-abierta-para-el-anuxe1lisis-de-sondeos}}

\textbf{Revista:} Geogaceta (Sociedad Geológica de España). Alternativa:
Spanish Journal of Palaeontology.

\hypertarget{tuxedtulo-provisional}{%
\subsubsection{Título provisional}\label{tuxedtulo-provisional}}

SedCoreTool: una herramienta abierta para cuantificar la bioturbación e
integrar datos multiproxy en sondeos {[}TODO{]}

\hypertarget{idea-1}{%
\subsubsection{Idea}\label{idea-1}}

Presentar el programa a la comunidad española con un ejemplo
reproducible sobre datos públicos (IODP U1385I-1H, margen ibérico) y el
proyecto de entrenamiento: imagen tratada $\rightarrow$ porcentaje de bioturbación,
BI e ii $\rightarrow$ columna compuesta con registros geoquímicos $\rightarrow$ familias de
intervalos y clasificación $\rightarrow$ mapa por edad.

\hypertarget{esquema}{%
\subsubsection{Esquema}\label{esquema}}

\begin{enumerate}
\def\labelenumi{\arabic{enumi}.}
\tightlist
\item
  Introducción: subjetividad de los índices visuales y falta de
  herramientas abiertas que integren imagen y datos.
\item
  La herramienta: flujo, formatos, reproducibilidad (proyecto guardado,
  semilla), aplicación web sin instalación.
\item
  Ejemplo con datos públicos (figura única con tres paneles).
\item
  Disponibilidad: código, DOI, licencia, aplicación en línea.
\end{enumerate}

\hypertarget{pendiente}{%
\subsubsection{Pendiente}\label{pendiente}}

\begin{itemize}
\tightlist
\item[$\square$]
  Normas de extensión y formato de la revista.
\item[$\square$]
  Figura y aplicación publicada en GitHub Pages.
\end{itemize}

\newpage

\hypertarget{revistas-candidatas}{%
\section{Revistas candidatas}\label{revistas-candidatas}}

No se incluyen factores de impacto ni APC: cambian cada año. Verificar
en la web de cada revista antes de enviar.

\hypertarget{software-y-muxe9todos-computacionales}{%
\subsection{Software y métodos
computacionales}\label{software-y-muxe9todos-computacionales}}

\begin{longtable}[]{@{}
  >{\raggedright\arraybackslash}p{(\columnwidth - 6\tabcolsep) * \real{0.2500}}
  >{\raggedright\arraybackslash}p{(\columnwidth - 6\tabcolsep) * \real{0.3056}}
  >{\raggedright\arraybackslash}p{(\columnwidth - 6\tabcolsep) * \real{0.2222}}
  >{\raggedright\arraybackslash}p{(\columnwidth - 6\tabcolsep) * \real{0.2222}}@{}}
\toprule\noalign{}
\begin{minipage}[b]{\linewidth}\raggedright
Revista
\end{minipage} & \begin{minipage}[b]{\linewidth}\raggedright
Editorial
\end{minipage} & \begin{minipage}[b]{\linewidth}\raggedright
Acceso
\end{minipage} & \begin{minipage}[b]{\linewidth}\raggedright
Encaje
\end{minipage} \\
\midrule\noalign{}
\endhead
\bottomrule\noalign{}
\endlastfoot
Journal of Open Source Software (JOSS) & Open Journals & Abierto, sin
coste & Artículo corto + revisión del código en GitHub. Ideal para
IchAnTool 3. \\
Computers \& Geosciences & Elsevier & Híbrido & Donde se publicó la v2.1
(2020). Bueno para un método nuevo (segmentación automática). \\
Applied Computing and Geosciences & Elsevier & Abierto & Revista hermana
de C\&G, orientada a aplicaciones y software. \\
Artificial Intelligence in Geosciences & KeAi & Abierto & Encaje directo
si el aporte principal es el modelo de deep learning. \\
Geoscientific Instrumentation, Methods and Data Systems & Copernicus &
Abierto & Métodos y sistemas de datos; alternativa a C\&G. \\
\end{longtable}

\hypertarget{datos}{%
\subsection{Datos}\label{datos}}

\begin{longtable}[]{@{}
  >{\raggedright\arraybackslash}p{(\columnwidth - 6\tabcolsep) * \real{0.2500}}
  >{\raggedright\arraybackslash}p{(\columnwidth - 6\tabcolsep) * \real{0.3056}}
  >{\raggedright\arraybackslash}p{(\columnwidth - 6\tabcolsep) * \real{0.2222}}
  >{\raggedright\arraybackslash}p{(\columnwidth - 6\tabcolsep) * \real{0.2222}}@{}}
\toprule\noalign{}
\begin{minipage}[b]{\linewidth}\raggedright
Revista
\end{minipage} & \begin{minipage}[b]{\linewidth}\raggedright
Editorial
\end{minipage} & \begin{minipage}[b]{\linewidth}\raggedright
Acceso
\end{minipage} & \begin{minipage}[b]{\linewidth}\raggedright
Encaje
\end{minipage} \\
\midrule\noalign{}
\endhead
\bottomrule\noalign{}
\endlastfoot
Earth System Science Data (ESSD) & Copernicus & Abierto & Dataset
anotado de imágenes de testigos y máscaras de bioturbación. \\
Scientific Data & Springer Nature & Abierto & Data descriptor; alta
visibilidad fuera de geología. \\
Geoscience Data Journal & Wiley & Abierto & Alternativa para datasets
geocientíficos. \\
\end{longtable}

\hypertarget{icnologuxeda-sedimentologuxeda-y-paleoambiente}{%
\subsection{Icnología, sedimentología y
paleoambiente}\label{icnologuxeda-sedimentologuxeda-y-paleoambiente}}

\begin{longtable}[]{@{}
  >{\raggedright\arraybackslash}p{(\columnwidth - 6\tabcolsep) * \real{0.2500}}
  >{\raggedright\arraybackslash}p{(\columnwidth - 6\tabcolsep) * \real{0.3056}}
  >{\raggedright\arraybackslash}p{(\columnwidth - 6\tabcolsep) * \real{0.2222}}
  >{\raggedright\arraybackslash}p{(\columnwidth - 6\tabcolsep) * \real{0.2222}}@{}}
\toprule\noalign{}
\begin{minipage}[b]{\linewidth}\raggedright
Revista
\end{minipage} & \begin{minipage}[b]{\linewidth}\raggedright
Editorial
\end{minipage} & \begin{minipage}[b]{\linewidth}\raggedright
Acceso
\end{minipage} & \begin{minipage}[b]{\linewidth}\raggedright
Encaje
\end{minipage} \\
\midrule\noalign{}
\endhead
\bottomrule\noalign{}
\endlastfoot
Ichnos & Taylor \& Francis & Híbrido & Revista de referencia en
icnología; validación frente a BI/ii visual. \\
Palaeogeography, Palaeoclimatology, Palaeoecology & Elsevier & Híbrido &
Aplicación paleoambiental (p.~ej. Último Ciclo Glacial en el
Mediterráneo). \\
Marine Geology & Elsevier & Híbrido & Aplicación en testigos marinos
(IODP/ODP). \\
Sedimentary Geology & Elsevier & Híbrido & Enfoque
sedimentológico/icnofábrica. \\
Palaios & SEPM & Híbrido & Paleoecología e icnología. \\
Lethaia & Scandinavian University Press & Abierto & Paleontología e
icnología. \\
Marine and Petroleum Geology / AAPG Bulletin & Elsevier / AAPG & Híbrido
& Si se orienta a reservorios y caracterización de testigos en
industria. \\
\end{longtable}

\hypertarget{divulgaciuxf3n-tuxe9cnica-en-espauxf1ol}{%
\subsection{Divulgación técnica en
español}\label{divulgaciuxf3n-tuxe9cnica-en-espauxf1ol}}

\begin{longtable}[]{@{}
  >{\raggedright\arraybackslash}p{(\columnwidth - 4\tabcolsep) * \real{0.3214}}
  >{\raggedright\arraybackslash}p{(\columnwidth - 4\tabcolsep) * \real{0.3929}}
  >{\raggedright\arraybackslash}p{(\columnwidth - 4\tabcolsep) * \real{0.2857}}@{}}
\toprule\noalign{}
\begin{minipage}[b]{\linewidth}\raggedright
Revista
\end{minipage} & \begin{minipage}[b]{\linewidth}\raggedright
Editorial
\end{minipage} & \begin{minipage}[b]{\linewidth}\raggedright
Encaje
\end{minipage} \\
\midrule\noalign{}
\endhead
\bottomrule\noalign{}
\endlastfoot
Geogaceta & Sociedad Geológica de España & Nota corta presentando la
herramienta a la comunidad española. \\
Spanish Journal of Palaeontology & Sociedad Española de Paleontología &
Aplicaciones icnológicas. \\
\end{longtable}

\hypertarget{criterios-para-elegir}{%
\subsection{Criterios para elegir}\label{criterios-para-elegir}}

\begin{enumerate}
\def\labelenumi{\arabic{enumi}.}
\tightlist
\item
  ¿El aporte es el software, los datos, el método o la interpretación
  geológica? Eso decide la familia de revistas.
\item
  Acceso abierto: aumenta adopción de la herramienta; JOSS y Copernicus
  no cobran o cobran poco.
\item
  Plazo: JOSS y las revistas de datos suelen ser más rápidas que las de
  paleoambiente.
\end{enumerate}

\hypertarget{refs}{}
\begin{CSLReferences}{1}{0}
\leavevmode\vadjust pre{\hypertarget{ref-casanova2020}{}}%
Casanova-Arenillas, Santiago, Francisco J. Rodríguez-Tovar, y Francisca
Martínez-Ruiz. 2020. {«Applied ichnology in sedimentary geology: Python
scripts as a method to automatize ichnofabric analysis in marine core
images»}. \emph{Computers \& Geosciences} 136: 104407.
\url{https://doi.org/10.1016/j.cageo.2020.104407}.

\leavevmode\vadjust pre{\hypertarget{ref-dorador2014}{}}%
Dorador, Javier, y Francisco J. Rodríguez-Tovar. 2014. {«Quantitative
estimation of bioturbation based on digital image analysis»}.
\emph{Marine Geology} 349: 55-60.

\leavevmode\vadjust pre{\hypertarget{ref-dorador2018}{}}%
---------. 2018. {«High-resolution image treatment in ichnological core
analysis: Initial steps, advances and prospects»}. \emph{Earth-Science
Reviews} 177: 226-37.

\leavevmode\vadjust pre{\hypertarget{ref-droser1986}{}}%
Droser, Mary L., y David J. Bottjer. 1986. {«A semiquantitative field
classification of ichnofabric»}. \emph{Journal of Sedimentary Petrology}
56: 558-59.

\leavevmode\vadjust pre{\hypertarget{ref-reineck1963}{}}%
Reineck, H.-E. 1963. {«Sedimentgef{ü}ge im Bereich der s{ü}dlichen
Nordsee»}. \emph{Abhandlungen der Senckenbergischen Naturforschenden
Gesellschaft} 505: 1-138.

\leavevmode\vadjust pre{\hypertarget{ref-taylor1993}{}}%
Taylor, A. M., y R. Goldring. 1993. {«Description and analysis of
bioturbation and ichnofabric»}. \emph{Journal of the Geological Society}
150: 141-48.

\end{CSLReferences}

\end{document}
